% Patch 0620 A5_F^(3): Third-Order Closed-Form Substrate-Current Coefficients
% at Face-Aligned Reading C in the 600-Cell
%
% OPEN-FP-F1-5 Sequence-3B closure artifact (THEO-DSL-11 candidate). Computes
% the O(delta^3) face-aligned coefficients alpha_3^(rho), alpha_3^(ax) in the
% 2D V_4-invariant subspace span{n_rho, n_ax} established at THEO-DSL-9 (Patch
% 0615) -- the V_4 correction to THEO-DSL-8 multi-AI confirmed at Patch 0617.
% Sketch-document Layer 3 under (H5_F^(3)) -- the natural extension of (H5_F)'s
% W(P)=1 ansatz from P_2(v_host) to P_3(v_host), per DG-F15A-4 default-(A) per-
% variant scope qualifier and DG-F15A-5 default-(A) single-theorem registration.
%
% CHANGELOG
%   v1.0 (Patch 0620, Session 147, 28 May 2026): higher-order extension of
%     THEO-DSL-9 (face k=2, Patch 0615 + Patch 0617 multi-AI confirmation).
%   v1.1 (Patch 0629, Session 147, 28 May 2026): consolidated wording-fix
%     revision pass per multi-AI review backlog from Patch 0622 (ChatGPT +
%     Copilot reviews; ChatGPT/Grok/Copilot all CONFIRMED). Specific changes:
%       - ChatGPT 1: §Robustness wording correction. The phrase "rotates
%         n_Fperp within T_{v_host} S^3" implicitly treated n_Fperp as
%         tangential; it has a radial component along v_host. Replaced
%         with explicit "rotates the tangential axial component while
%         preserving the same radial projection along v_host" framing.
%       - ChatGPT 2: §Robustness softening: the adapted 2D subspace
%         span{n_rho, n_ax} is face-dependent (n_ax depends on {u_i, u_j});
%         clarified that the coefficient pair is orbit-invariant in the
%         CORRESPONDING ADAPTED FRAME for each face choice.
%       - ChatGPT 3: Explicit denominator-simplification note for k=3:
%         Theorem 1 predicts ambient sqrt(3)*Q[phi]/9; observed denominator
%         is 3 (which divides 9). Same pattern as THEO-DSL-12 at k=4.
%         Added explicit text noting ambient-vs-observed denominator
%         structure to avoid denominator nitpicks.
%       - Copilot 1: Parity Theorem labelled "the Parity Theorem"
%         explicitly for visible cross-reference; the theorem already has
%         a label (thm:parity-structure) and a descriptive name in its
%         environment header, but Copilot wanted higher visibility.
%       - Copilot 2: Anti-erasure remark expanded with a short proof
%         sketch that no Q[phi]-basis change of span{n_rho, n_ax} can
%         absorb the sqrt(3) without reintroducing it elsewhere
%         (the sqrt(3) lives in the path-product weight, not in the
%         frame; any basis substitution preserves the path-product
%         structure).
%       - Copilot 3: Added a short geometric remark explaining why
%         |v_host + u_i + u_j|^2 = 3 + 3*phi = 3*phi^2 -- specifically,
%         the dot-product structure of three mutually adjacent 600-cell
%         vertices on a triangular face.
%     No theorem changes; no claim weakened or strengthened. All wording
%     fixes preserve the artifact's mathematical content; revisions are
%     exposition-level only.
%
% ** SUBSTANTIVE NEW FINDING -- the central content of this artifact **:
%
% The face-aligned k=3 coefficients are NOT in Q[phi]; they lie in sqrt(3)*Q[phi]/3:
%   alpha_3^(rho, face) = (87 - 53 phi) * sqrt(3) / 3  ~ +0.71834
%   alpha_3^(ax)        = (41 - 28 phi) * sqrt(3) / 3  ~ -2.48547
%
% Algebraic origin: |v_host + u_i + u_j|^2 = 3 + 3 phi = 3 phi^2, so
%   |v_host + u_i + u_j| = phi * sqrt(3),
% and n_Fperp = (v_host + u_i + u_j) / (phi*sqrt(3)) carries an unavoidable
% 1/sqrt(3) factor. Each factor (e_j . n_Fperp) carries 1/sqrt(3); the k-fold
% product has (1/sqrt(3))^k. At EVEN k the sqrt(3) factors pair up to a
% rational; at ODD k a single sqrt(3) survives. Hence:
%   - vertex variant (all k): clean Q[phi]              -- n_rho = v_host has no sqrt(3)
%   - edge variant (all k):   clean Q[phi]              -- n_edge = phi*(u_1-v_h) has no sqrt(3)
%   - face variant, k even:   clean Q[phi]              -- THEO-DSL-9 confirms at k=2
%   - face variant, k odd:    sqrt(3) * Q[phi] / 3^((k-1)/2 + 1)
%                                                       -- this artifact at k=3
%
% This is the **first F.1 result outside Q[phi]** in the entire programme. The
% parity-dependent algebraic structure -- Q[phi] vs sqrt(3)*Q[phi] depending on
% the parity of k at face-aligned -- is registered here as a substantive new
% feature, NOT a computational artifact (verified at 80-digit mpmath precision
% + PSLQ failure in {1, phi} basis + PSLQ success in {1, phi, sqrt(3), sqrt(3)*phi}
% basis with the rational and Q[phi] coefficients exactly zero, leaving only the
% sqrt(3) and sqrt(3)*phi components nonzero).
%
% The V_4-forced VANISHING of the two non-invariant components (Lemmas 1-2 of
% THEO-DSL-9 / Patch 0615, multi-AI confirmed at Patch 0617) persists at k=3:
%   alpha_3^(perp3) = 0  (forced by sigma_F in V_4, all orders)
%   alpha_3^(diff)  = 0  (forced by sigma_E in V_4, all orders)
%
% Mandatory vertex cross-check at k=3: alpha_3^(vertex) = -126 + 81 phi
% = 9(4-phi)/phi^3 ~ +5.061 (clean Q[phi]; reproduces THEO-DSL-10 cross-check).
% Path enumeration: 1728 = 12^3 directed 3-edge paths across 13 non-empty
% (s_v', s_v'') cells in the full 9-shell vertex classification (same reach as
% THEO-DSL-10 edge k=3, including B_5).
%
% Verify script: code/verify_face_alpha3_closure.py (ALL 30 PASS checks).
% Reasoning fragment: hardened_theorems/reasoning/0620.md.
% Version 1.2 --- 17 August 2026 (Patch 3214):
%   added a How-we-review methodology note stating plainly that AI panel
%   review is a self-applied verification method and not journal peer
%   review; twelve chirality papers retitled so the descriptive title
%   leads and the identifier is demoted to a subtitle. No claim,
%   equation, number or section changed.
% Version 1.1 --- 17 August 2026 (Patch 3213):
%   extended the generated identifier appendix to all five identifier
%   families (THEO-, FI-, PRED-, CONJ-, PH- as well as OPEN-), and
%   replaced review-convergence scores with AI review panel wording
%   where present. No claim, equation, number or section changed.
% Version 1.0 --- 17 August 2026 (Patch 3212):
%   added the generated plain-language appendix glossing the internal
%   programme identifiers used in this paper (founder jargon ruling, 16
%   Aug 2026). No claim, equation, number or section changed.
%   This is the FIRST version stamp assigned to this paper; 1.0 denotes
%   the first stamped revision, NOT a claim of v1.0 shipped grade.

\documentclass[11pt]{article}
\usepackage[utf8]{inputenc}
\usepackage{amsmath,amssymb,amsthm}
\usepackage[margin=1in]{geometry}
\usepackage{hyperref}
\usepackage{booktabs}

\theoremstyle{plain}
\newtheorem{theorem}{Theorem}
\newtheorem{lemma}{Lemma}
\newtheorem{corollary}{Corollary}
\theoremstyle{definition}
\newtheorem*{remark}{Remark}
\newtheorem*{antierasure}{Anti-erasure Remark}

\newcommand{\nhat}{\hat{n}}
\newcommand{\nrho}{\hat{n}_\rho}
\newcommand{\nax}{\hat{n}_{\text{ax}}}
\newcommand{\nperpF}{\hat{n}_{F\perp}}
\newcommand{\nperpthree}{\hat{n}_{\perp 3}}
\newcommand{\ndiff}{\hat{n}_{\text{diff}}}
\newcommand{\vhost}{v_{\text{host}}}
\newcommand{\ui}{u_i}
\newcommand{\uj}{u_j}
\newcommand{\jvec}{\vec{j}}
\newcommand{\Reals}{\mathbb{R}}
\newcommand{\Vfour}{V_4}

\title{Third-Order Closed-Form Substrate-Current Coefficients\\
at Face-Aligned Reading C in the 600-Cell\\
{\large (Sketch-Document Layer 3 under $(H5_F^{(3)})$; THEO-DSL-11 candidate;\\
OPEN-FP-F1-5 Sequence-3B; Patch 0620)\\[4pt]
{\normalsize \textbf{First F.1 result outside $\mathbb{Q}[\phi]$ --- odd-$k$ centroid-direction normalization signature}}}}
\author{CPP Programme}
\date{28 May 2026 (Session 147)\\[2pt]
{\small Version~1.0 --- 17 August 2026 (Patch 3212: added the generated plain-language appendix glossing the internal programme identifiers used in this paper (founder jargon ruling, 16 Aug 2026). No claim, equation, number or section changed.)}\\[2pt]
{\small Version~1.1 --- 17 August 2026 (Patch 3213: extended the generated identifier appendix to all five identifier families (THEO-, FI-, PRED-, CONJ-, PH- as well as OPEN-), and replaced review-convergence scores with AI review panel wording where present. No claim, equation, number or section changed.)}\\[2pt]
{\small Version~1.2 --- 17 August 2026 (Patch 3214: added a How-we-review methodology note stating plainly that AI panel review is a self-applied verification method and not journal peer review; twelve chirality papers retitled so the descriptive title leads and the identifier is demoted to a subtitle. No claim, equation, number or section changed.)}}

\begin{document}
\maketitle

\begin{abstract}
This artifact closes Sequence-3B of OPEN-FP-F1-5, extending THEO-DSL-9 (the
face-aligned $\mathcal{O}(\delta^2)$ closure with the $V_4$/2D structural
correction to THEO-DSL-8, multi-AI confirmed at Patch 0617) to the next order
in $\delta$. Under the $V_4$-invariant 2-dimensional subspace
$\mathrm{span}\{\nrho, \nax\}$ established at THEO-DSL-9, the order-three
substrate-current vector at face-aligned Reading C is given in closed form by
\begin{equation*}
\boxed{\;\;
  \jvec_3(\vhost) \;=\; \alpha_3^{(\rho)} \nrho \;+\; \alpha_3^{(\text{ax})} \nax,
  \quad
  \alpha_3^{(\rho)} = \frac{(87 - 53\phi)\sqrt{3}}{3},
  \quad
  \alpha_3^{(\text{ax})} = \frac{(41 - 28\phi)\sqrt{3}}{3}.
\;\;}
\end{equation*}
Numerically $\alpha_3^{(\rho)} \approx +0.71834$ and $\alpha_3^{(\text{ax})}
\approx -2.48547$.

\paragraph{Substantive new finding (the central content of this artifact).}
Unlike every prior closed-form coefficient in the F.1 programme --- which has
lived in $\mathbb{Q}[\phi]$ since the inception of the dynamical substrate law
arc --- the face-aligned $k=3$ coefficients live in $\sqrt{3}\cdot\mathbb{Q}[\phi]/3$.
The $\sqrt{3}$ factor is traceable to the exact-norm identity
$\|\vhost + \ui + \uj\| = \phi\sqrt{3}$, which forces an unavoidable $1/\sqrt{3}$
into the face-centroid direction $\nperpF$. At even $k$, pairs of factors
$1/\sqrt{3} \cdot 1/\sqrt{3} = 1/3$ cancel back to rationals (THEO-DSL-9 at
$k=2$, clean $\mathbb{Q}[\phi]$); at odd $k$ a single $\sqrt{3}$ survives. The
parity-dependent algebraic structure --- $\mathbb{Q}[\phi]$ at even $k$ versus
$\sqrt{3}\cdot\mathbb{Q}[\phi]$ at odd $k$ for the face-aligned variant ---
is a real feature of the geometry, not a computational artifact. \textbf{This
is the first F.1 result outside $\mathbb{Q}[\phi]$.}

\paragraph{Structural inheritance preserved.} The $V_4$-forced vanishing of
the two non-invariant components carries over to $k=3$ unconditionally from
THEO-DSL-9 Lemmas 1--2 (multi-AI confirmed at Patch 0617):
$\alpha_3^{(\perp 3)} = 0$ (forced by $\sigma_F$) and $\alpha_3^{(\text{diff})}
= 0$ (forced by $\sigma_E$). The vertex-aligned cross-check at $k=3$
reproduces $\alpha_3^{(\text{vertex})} = -126 + 81\phi = 9(4 - \phi)/\phi^3
\approx +5.061$ (matching the THEO-DSL-10 cross-check, clean $\mathbb{Q}[\phi]$).
The result is verified by independent numerical assembly of all
$1728 = 12^3$ directed 3-edge paths from the host vertex, decomposing in the
orthonormal basis $\{\nrho, \nax, \nperpthree, \ndiff\}$ over 13 non-empty
$(s_{v'}, s_{v''})$ cells with the path reach extending to shell $B_5$ at
this order. The full computation is sketch-document Layer 3 under the natural
extension $(H5_F^{(3)})$ of the path-class weight ansatz $W(P) = 1$ from
$\mathcal{P}_2(\vhost)$ to $\mathcal{P}_3(\vhost)$.
\end{abstract}

\section{Scope and provenance}\label{sec:scope}

This artifact registers a new candidate hardened theorem THEO-DSL-11:
the third-order ($\mathcal{O}(\delta^3)$) closed-form substrate-current
coefficients at face-aligned Reading C in the 600-cell, under the
2-dimensional $V_4$-invariant subspace established at THEO-DSL-9 (Patch
0615), with the $V_4$ structural correction to THEO-DSL-8 multi-AI confirmed
at Patch 0617. It is the direct higher-order continuation of THEO-DSL-9 and
the face-aligned cousin of THEO-DSL-10 (edge $k=3$, Patch 0618).

\paragraph{Inheritance from prior theorems.}
\begin{itemize}
\item \textbf{THEO-DSL-9} (candidate, Patch 0615; multi-AI confirmed Patch 0617):
  face-aligned $\mathcal{O}(\delta^2)$ closure $\alpha_2^{(\rho)} = -14 + 7\phi
  = -7/\phi^2$ and $\alpha_2^{(\text{ax})} = -6 + 2\phi$, in clean $\mathbb{Q}[\phi]$;
  $V_4$/2D structural sub-result (Lemmas 1--2; cross-reviewer THREE-FOR-THREE
  no-refutation; $V_4$-saturation proven by full 120-element host-stabilizer
  enumeration). The structural lemmas carry over unconditionally to $k=3$.
\item \textbf{THEO-DSL-10} (candidate, Patch 0618): edge-aligned $\mathcal{O}(\delta^3)$
  closure $\alpha_3^{(\rho)} = -84 + 54\phi$ and $\alpha_3^{(\text{edge})} =
  -63/2 + (21/2)\phi$, in clean $\mathbb{Q}[\phi]$. The path-enumeration
  machinery (1728 paths, full 9-shell classification) is reused here at
  face-aligned.
\item \textbf{THEO-DSL-5} (Patch 0591), \textbf{THEO-DSL-7} (Patch 0613): vertex
  and edge $k=2$ closures, $\mathbb{Q}[\phi]$. Used as additional cross-checks
  via the same assembly machinery.
\end{itemize}

\paragraph{Hypotheses retained.}
$(H1_F)$ Reading C with host $\vhost$ and face-centroid anchor
$\nperpF = (\vhost + \ui + \uj)/\|\vhost + \ui + \uj\|$ for a chosen triangular
face $\{\vhost, \ui, \uj\}$ at $\vhost$ ($\ui, \uj \in S_1$ with
$\ui \cdot \uj = \phi/2$);
$(H2)$ Mechanism A (anisotropic-rate path-integral first-order substrate-current
ansatz with rate $r(\hat e) = r_0(1 + \delta \hat e \cdot \nperpF)$);
$(H3_F)$ residual symmetry $\Vfour = \{e, \sigma_E, \sigma_F, C_2\}$ at $\vhost$,
saturating the face stabilizer ($|I_h|/|\text{face orbit}| = 120/30 = 4$),
multi-AI confirmed at Patch 0617;
$(H4)$ the 600-cell unit-circumradius edge graph;
$(H5_F^{(3)})$ the path-class weight $W(P) = 1$ extended from
$\mathcal{P}_2(\vhost)$ to $\mathcal{P}_3(\vhost)$ (no $k$-dependent
reweighting).

The result is sketch-document Layer 3 under $(H5_F^{(3)})$: the assembly is
unconditional given the ansatz, and the closed forms below are exact in
$\sqrt{3} \cdot \mathbb{Q}[\phi]/3$ rather than numerical. The structural
sub-result ($\jvec_3^{\text{face}} \in \mathrm{span}\{\nrho, \nax\}$) inherited
from THEO-DSL-9 is publication-grade Layer 3 unconditional (it depends only on
$V_4$-equivariance, not on any path-class weight).

\section{The face frame and the $\sqrt{3}$ identity}\label{sec:sqrt3}

The orthonormal frame at the chosen face $\{\vhost, \ui, \uj\}$:
\begin{align}
\nrho       &\;=\; \vhost                              & &\text{(radial)}\\
\nax        &\;=\; \ui + \uj - \phi\,\vhost            & &\text{(face axial; \textbf{exact unit norm}, see Lemma~\ref{lem:nax-norm})}\\
\nperpthree &\;=\; \text{unit vector} \perp \mathrm{span}\{\vhost, \ui, \uj\} & &\text{(out-of-3-flat)}\\
\ndiff      &\;=\; (\ui - \uj)/\|\ui - \uj\|           & &\text{(in-face anti-symmetric)}
\end{align}
and the face-centroid direction
\begin{equation}\label{eq:nFperp-def}
\nperpF \;=\; \frac{\vhost + \ui + \uj}{\|\vhost + \ui + \uj\|}.
\end{equation}

\begin{lemma}[exact $\nax$ unit norm]\label{lem:nax-norm}
$\|\ui + \uj - \phi\,\vhost\|^2 = 1$ exactly.
\end{lemma}
\begin{proof}
Direct expansion:
\begin{align*}
\|\ui + \uj - \phi\vhost\|^2
   &= |\ui|^2 + |\uj|^2 + \phi^2 + 2(\ui \cdot \uj) - 2\phi(\ui \cdot \vhost) - 2\phi(\uj \cdot \vhost)\\
   &= 1 + 1 + \phi^2 + 2\cdot\tfrac{\phi}{2} - 2\phi\cdot\tfrac{\phi}{2} - 2\phi\cdot\tfrac{\phi}{2}\\
   &= 2 + \phi^2 + \phi - 2\phi^2
   = 2 + \phi - \phi^2
   = 2 + \phi - (\phi + 1) = 1. \qedhere
\end{align*}
\end{proof}

\begin{lemma}[the $\sqrt{3}$ identity]\label{lem:phi-sqrt3}
$\|\vhost + \ui + \uj\| = \phi\sqrt{3}$ exactly.
\end{lemma}
\begin{proof}
\begin{align*}
\|\vhost + \ui + \uj\|^2
   &= 3 + 2(\vhost \cdot \ui) + 2(\vhost \cdot \uj) + 2(\ui \cdot \uj)\\
   &= 3 + 3 \cdot 2 \cdot \tfrac{\phi}{2} = 3 + 3\phi = 3(1+\phi) = 3\phi^2.
\end{align*}
Taking the positive square root, $\|\vhost + \ui + \uj\| = \phi\sqrt{3}$.
\end{proof}

\begin{corollary}[$\nperpF$ algebraic structure]\label{cor:nFperp-sqrt3}
$\nperpF \in \mathbb{Q}[\phi,\,\sqrt{3}\,]^4$, with the unique $\sqrt{3}$
appearing only in the denominator. Specifically,
\begin{equation}\label{eq:nFperp-explicit}
\nperpF \;=\; \frac{\vhost + \ui + \uj}{\phi\sqrt{3}}
       \;=\; \frac{1}{\sqrt{3}} \cdot \frac{\vhost + \ui + \uj}{\phi},
\end{equation}
so $\nperpF = \sqrt{3}^{-1} \cdot w$ for some $w \in \mathbb{Q}[\phi]^4$.
\end{corollary}

In contrast, the vertex and edge frame vectors are entirely in
$\mathbb{Q}[\phi]^4$:
$\nrho = \vhost \in \mathbb{Q}[\phi]^4$ trivially, and
$\hat n_{\text{edge}} = \phi(u_1 - \vhost) \in \mathbb{Q}[\phi]^4$ since the
edge length is $|u_1 - \vhost| = 1/\phi$. The face variant is unique in
carrying a $\sqrt{3}$ factor.

\section{Parity-dependent algebraic structure}\label{sec:parity}

The path-integral assembly at order $k$:
\begin{equation}\label{eq:jk-assembly}
\jvec_k(\vhost) \;=\; \sum_{P \in \mathcal{P}_k(\vhost)}
   \prod_{j=1}^{k} (\hat e_j \cdot \nperpF) \cdot \hat e_1,
\end{equation}
where the assembly is over directed $k$-edge paths $\vhost \to u_1 \to \cdots
\to u_k$ and $\hat e_j$ is the unit edge direction of step $j$.

By Corollary~\ref{cor:nFperp-sqrt3}, each $\hat e_j \cdot \nperpF$ has the
form $c_j/\sqrt{3}$ where $c_j \in \mathbb{Q}[\phi]$. The $k$-fold product
gives
\begin{equation}\label{eq:k-fold-product}
\prod_{j=1}^{k} (\hat e_j \cdot \nperpF) \;=\; \frac{\prod_j c_j}{(\sqrt{3})^k}
   \;=\; \frac{\prod_j c_j}{3^{k/2}}
\end{equation}
which equals
\begin{equation*}
\begin{cases}
\dfrac{1}{3^{k/2}}\,\big(\text{element of } \mathbb{Q}[\phi]\big) & \text{if } k \text{ even},\\[4pt]
\dfrac{\sqrt{3}}{3^{(k+1)/2}}\,\big(\text{element of } \mathbb{Q}[\phi]\big) & \text{if } k \text{ odd}.
\end{cases}
\end{equation*}
Summing over $|\mathcal{P}_k(\vhost)| = 12^k$ paths and projecting onto a
basis vector in $\mathbb{Q}[\phi]^4$ (such as $\nrho$ or $\nax$) preserves
this structure. Hence:

\begin{theorem}[parity-dependent algebraic structure --- the \emph{Parity Theorem}]\label{thm:parity-structure}
At face-aligned Reading C in the 600-cell with the $\nperpF$ substrate
primitive direction, the $k$-th-order substrate-current coefficients
$\alpha_k^{(\rho, \text{face})}$ and $\alpha_k^{(\text{ax})}$ (in the
$V_4$-invariant 2D subspace $\mathrm{span}\{\nrho, \nax\}$) satisfy:
\begin{equation}
\alpha_k^{(\rho, \text{face})},\;\alpha_k^{(\text{ax})}
\;\in\;
\begin{cases}
\mathbb{Q}[\phi] / 3^{k/2}        & \text{if } k \text{ even},\\[4pt]
\sqrt{3}\cdot\mathbb{Q}[\phi] / 3^{(k+1)/2} & \text{if } k \text{ odd}.
\end{cases}
\end{equation}
The vertex (THEO-DSL-3, -4, -5; THEO-DSL-10 cross-check) and edge (THEO-DSL-7,
-10) variants live in $\mathbb{Q}[\phi]$ at all $k$, since their substrate
primitive directions ($\nrho = \vhost$, $\hat n_{\text{edge}} \propto u_1 - \vhost$)
carry no $\sqrt{3}$.

Hereafter this result is referenced as \emph{the Parity Theorem}.
\end{theorem}

\begin{proof}[Sketch]
The path-class weight $W(P) = 1$ ansatz $(H5_F^{(k)})$ + the rate-function
factorisation~\eqref{eq:jk-assembly} + the $\sqrt{3}$-only-in-$\nperpF$
structure of Corollary~\ref{cor:nFperp-sqrt3} give a sum of $12^k$
$\mathbb{Q}[\phi,\sqrt{3}]$-valued vectors of the form
$\sqrt{3}^{-k}\cdot(\text{element of }\mathbb{Q}[\phi]^4)$. Projection onto a
$\mathbb{Q}[\phi]^4$ basis vector yields a scalar of the form
$\sqrt{3}^{-k}\cdot(\text{rational of }\mathbb{Q}[\phi])$. The parity case
analysis on $\sqrt{3}^{-k}$ gives the stated dichotomy.
\end{proof}

\begin{remark}[ambient vs observed denominator]\label{rem:ambient-vs-observed}
At $k=3$ (odd), the Parity Theorem predicts the ambient denominator
$3^{(k+1)/2} = 3^2 = 9$; the observed coefficients in
Theorem~\ref{thm:dsl11} below have denominator 3 (which divides 9).
This is the same ``ambient prediction $\supseteq$ observed denominator''
pattern that recurs at $k=4$ in THEO-DSL-12, where Theorem~\ref{thm:parity-structure}
predicts ambient $\mathbb{Q}[\phi]/9$ and the observed coefficients have
denominators 2 and 3 (both dividing 9). The Parity Theorem gives an
\emph{upper bound} on the denominator at order $k$, not an equality;
the actual coefficient may have a tighter denominator if particular
sub-sums happen to cancel.
\end{remark}

\begin{remark}[geometric origin of the $\sqrt{3}$ factor]\label{rem:geometric-origin}
The $\sqrt{3}$ enters via $|\vhost + \ui + \uj|^2 = 3 + 3\phi = 3\phi^2$,
so $|\vhost + \ui + \uj| = \phi\sqrt{3}$, and the face-centroid direction
$\nperpF = (\vhost + \ui + \uj)/(\phi\sqrt{3})$ carries an unavoidable
$1/\sqrt{3}$ factor. The geometric reason $|\vhost + \ui + \uj|^2 = 3 + 3\phi$
is that $(\vhost, \ui, \uj)$ are three mutually adjacent 600-cell vertices
(a triangular face), each pair at inner product $\phi/2$:
\begin{align*}
|\vhost + \ui + \uj|^2
   &= |\vhost|^2 + |\ui|^2 + |\uj|^2
      + 2\langle\vhost,\ui\rangle + 2\langle\vhost,\uj\rangle + 2\langle\ui,\uj\rangle\\
   &= 3 + 3 \cdot 2 \cdot \tfrac{\phi}{2} = 3 + 3\phi = 3(1+\phi) = 3\phi^2.
\end{align*}
The squared norm $3 + 3\phi$ is \emph{not} a perfect square in $\mathbb{Q}[\phi]$
(since $\phi$ satisfies $\phi^2 = \phi + 1$, the only square roots in
$\mathbb{Q}[\phi]$ of $3 + 3\phi = 3\phi^2$ are $\pm\phi\sqrt{3}$ in
$\mathbb{Q}[\phi, \sqrt{3}]$, which extend $\mathbb{Q}[\phi]$ properly).
This is the algebraic-geometric reason the face-aligned variant escapes
$\mathbb{Q}[\phi]$ at odd $k$ where the $\sqrt{3}^{-k}$ factor survives parity
cancellation.
\end{remark}

\begin{remark}
This is the first F.1 result outside $\mathbb{Q}[\phi]$ in the dynamical
substrate law programme. It is also the first case in which the algebraic
type of the closed-form coefficient depends non-trivially on the parity of
the order $k$. The structure is a genuine geometric feature of the face-
centroid direction in the 600-cell --- specifically, the fact that the
sum-of-three-unit-vectors $\vhost + \ui + \uj$ has squared norm
$3 + 3\phi = 3\phi^2$ rather than a perfect square in $\mathbb{Q}[\phi]$
(Remark~\ref{rem:geometric-origin}) --- not a computational artifact.
\end{remark}

\section{Closed forms at $k=3$}\label{sec:results}

By Theorem~\ref{thm:parity-structure}, the face-aligned $k=3$ coefficients
must lie in $\sqrt{3}\cdot\mathbb{Q}[\phi]/3$. The explicit values:

\begin{theorem}[THEO-DSL-11 candidate: face-aligned $\mathcal{O}(\delta^3)$
coefficients]\label{thm:dsl11}
Under hypotheses $(H1_F)$--$(H4)$ and $(H5_F^{(3)})$,
\begin{align}
\jvec_3(\vhost) &= \alpha_3^{(\rho)} \nrho \;+\; \alpha_3^{(\text{ax})} \nax, \\
\alpha_3^{(\rho)} &= \frac{(87 - 53\phi)\sqrt{3}}{3}
                   = \sqrt{3}\left(29 - \frac{53\phi}{3}\right)
                   \approx +0.71834, \\
\alpha_3^{(\text{ax})} &= \frac{(41 - 28\phi)\sqrt{3}}{3}
                   = \sqrt{3}\left(\frac{41 - 28\phi}{3}\right)
                   \approx -2.48547.
\end{align}
Both coefficients are independent of the choice of incident face
$\{\ui, \uj\}$ (verified explicitly over all 30 triangular faces incident to
$\vhost$, \S\ref{sec:robustness}). The two non-invariant components vanish
identically:
\begin{equation}
\alpha_3^{(\perp 3)} = 0, \quad \alpha_3^{(\text{diff})} = 0,
\end{equation}
forced unconditionally by $V_4$-equivariance (Lemmas 1--2 of THEO-DSL-9 /
Patch 0615; multi-AI confirmed at Patch 0617).
\end{theorem}

\begin{proof}[Computation]
The vector~\eqref{eq:jk-assembly} is evaluated by enumerating all
$1728 = 12^3$ directed 3-edge paths from $\vhost$ and summing the product of
three dot products against $\hat e_1$ (the first-edge direction). The
$V_4$-forced direction of the result --- $\jvec_3 \in \mathrm{span}\{\nrho,
\nax\}$ --- is the same structural lemma used at $k = 2$ (THEO-DSL-9 \S2.1)
and carries over to $k=3$ without modification (the lemma depends only on
$V_4$-equivariance of the assembly and $V_4$-fixedness of $\nperpF$). The
non-zero components in the basis $\{\nrho, \nax\}$ are extracted via
inner products in the exact-unit-norm frame (Lemma~\ref{lem:nax-norm}); the
result is identified with the closed forms above via PSLQ at 80-digit
mpmath precision in the basis $\{1, \phi, \sqrt{3}, \sqrt{3}\phi\}$,
showing that the rational and $\mathbb{Q}[\phi]$-only components are
exactly zero while the $\sqrt{3}$ and $\sqrt{3}\phi$ components are
$(29, -53/3)$ and $(41/3, -28/3)$ respectively. The closed forms are then
verified to machine-double precision in the verify script
\texttt{code/verify\_face\_alpha3\_closure.py} (all 30 PASS checks).
\end{proof}

\begin{corollary}[$k=3$ vertex cross-check]
The same 1728-path assembly with $\nhat = \vhost$ (vertex-aligned cross-check)
reproduces $\alpha_3^{(\text{vertex})} = -126 + 81\phi = 9(9\phi - 14) = 9(4
- \phi)/\phi^3 \approx +5.061$, in clean $\mathbb{Q}[\phi]$ as predicted by
THEO-DSL-4 + Theorem~\ref{thm:parity-structure}. The $k=1, k=2$ cross-checks
yield $3/\phi^2 = 6 - 3\phi$ and $-9/\phi^2 = -18 + 9\phi$ respectively,
locking the assembly convention to the one used in THEO-DSL-5/-7/-9/-10.
\end{corollary}

\section{Per-2D-shell-class decomposition}\label{sec:per-class}

The 1728 directed 3-edge paths distribute across the same 13 non-empty
$(s_{v'}, s_{v''})$ cells as the edge $k=3$ case (THEO-DSL-10); the cell
structure reflects the geometry of the 600-cell rather than the choice of
substrate direction, so the cell counts and shell reach are identical. The
distribution into the per-cell coefficient contributions, however, depends
on $\nperpF$ and is distinct from the edge case. The per-class numerics
(decimal approximations in $\sqrt{3}\cdot\mathbb{Q}[\phi]/3$):

\begin{center}
\small
\begin{tabular}{l c r r}
\toprule
$(s_{v'}, s_{v''})$ & $\#P$ & $\alpha_3^{(\rho)}$-cell & $\alpha_3^{(\text{ax})}$-cell \\
\midrule
$(B_0, B_1)$    & $144$  & $-1.5636$  & $-2.5820$ \\
$(B_1, B_0)$    & $60$   & $\phantom{-}0.0000$  & $+0.2887$ \\
$(B_1, B_1)$    & $300$  & $-0.3139$  & $-0.0240$ \\
$(B_1, B_2)$    & $300$  & $-0.8219$  & $-2.3982$ \\
$(B_1, B_3)$    & $60$   & $-0.2272$  & $-0.7731$ \\
$(B_2, B_1)$    & $180$  & $-4.4954$  & $-6.3668$ \\
$(B_2, B_2)$    & $180$  & $-0.4543$  & $-0.9375$ \\
$(B_2, B_3)$    & $180$  & $+2.0432$  & $+2.4180$ \\
$(B_2, B_4)$    & $180$  & $+6.0843$  & $+7.8473$ \\
$(B_3, B_1)$    & $12$   & $-0.7601$  & $-1.1196$ \\
$(B_3, B_2)$    & $60$   & $-1.5921$  & $-2.4938$ \\
$(B_3, B_4)$    & $60$   & $+1.9813$  & $+2.5289$ \\
$(B_3, B_5)$    & $12$   & $+0.8379$  & $+1.1266$ \\
\midrule
\textbf{Total}  & $1728$ & $+0.7183 \;=\; (87 - 53\phi)\sqrt{3}/3$
                         & $-2.4855 \;=\; (41 - 28\phi)\sqrt{3}/3$\\
\bottomrule
\end{tabular}
\end{center}

\paragraph{Observations.}
\begin{enumerate}
\item All 13 cells have $\hat n_{\perp 3}$-component and $\hat n_{\text{diff}}$-
  component $\sim 10^{-16}$ (machine zero), confirming the $V_4$-forced
  vanishing cell-by-cell, not merely in the total --- a stronger structural
  fact than the global zero.
\item The two largest single contributors are $(B_2, B_4)$ (paths reaching the
  equator at $v''$) with $+6.08, +7.85$ and $(B_2, B_1)$ with $-4.50, -6.37$,
  partially cancelling.
\item The path reach to $B_5$ in the $(B_3, B_5)$ cell (12 paths) is small
  but non-vanishing: $+0.84, +1.13$, signed differently from the edge case.
\end{enumerate}

\section{Robustness over face choice}\label{sec:robustness}

By THEO-DSL-9, the coefficient pair $(\alpha_3^{(\rho)},
\alpha_3^{(\text{ax})})$ is invariant under the $H_4$-orbit action on the 30
incident faces $\{\ui, \uj\}$ at $\vhost$. More precisely: choosing a
different face rotates the tangential axial component of $\nperpF$ while
preserving the same radial projection along $\vhost = \nrho$; correspondingly,
the adapted 2D subspace $\mathrm{span}\{\nrho, \nax\}$ \emph{itself} changes
with the chosen face (because $\nax = \ui + \uj - \phi\,\vhost$ depends
explicitly on $\{\ui, \uj\}$), but the coefficient pair
$(\alpha_3^{(\rho)}, \alpha_3^{(\text{ax})})$ is orbit-invariant in the
\emph{corresponding} adapted frame for each face choice. The verify
script confirms this explicitly: across all 30 faces, the standard
deviation of each coefficient over the orbit is $< 10^{-11}$ (double
precision); the $\hat n_{\perp 3}$ and $\hat n_{\text{diff}}$ components are
identically zero ($|\text{max}| < 10^{-12}$).

\section{Exclusions retained}\label{sec:exclusions}
\begin{enumerate}
\item[$(E1)$] No path-class weight other than $W(P) = 1$ at $k=3$ is
  considered; $(H5_F^{(3)})$ is the natural extension of $(H5_F)$ and is the
  same ansatz used at THEO-DSL-9.
\item[$(E2)$] Vertex (THEO-DSL-3,-4,-5; THEO-DSL-10 cross-check) and edge
  (THEO-DSL-7, -10) variants are separate artifacts at the same order.
\item[$(E3)$] No promotion of Layer 3 status to publication-grade: the
  underlying ansatz $(H5_F^{(3)})$ is sketch-document Layer 3, not
  axiomatically derived. The $V_4$/2D structural sub-results (vanishing of
  $\alpha_3^{(\perp 3)}$ and $\alpha_3^{(\text{diff})}$) are publication-grade
  Layer 3 unconditional, multi-AI confirmed at Patch 0617 in the THEO-DSL-9
  review cycle.
\item[$(E4)$] No claim that the parity-dependent $\mathbb{Q}[\phi]$ vs
  $\sqrt{3}\cdot\mathbb{Q}[\phi]$ structure (Theorem~\ref{thm:parity-structure})
  generalises to other Reading-C variants or other path-class weight ansatze.
  The $\sqrt{3}$ in $\nperpF$ is specific to the face-centroid direction at
  vertex-host convention.
\item[$(E5)$] No extension to $k \geq 4$ in this artifact (Sequence-4B and
  beyond are future OPEN-FP-F1-5 targets; $k=4$ face is predicted to be back
  in clean $\mathbb{Q}[\phi]/9$ by Theorem~\ref{thm:parity-structure}, but
  this prediction is not yet computationally verified).
\end{enumerate}

\section{Provenance and verification}\label{sec:provenance}

The exact closed forms above are obtained by:
\begin{enumerate}
\item Building the 600-cell from first principles (16-cell + tesseract +
  snub 24-cell with even-permutation parity) and verifying its 120-vertex,
  720-edge, 12-regular structure.
\item Verifying the exact-norm identities $\|\nax\| = 1$ (Lemma~\ref{lem:nax-norm})
  and $\|\vhost + \ui + \uj\| = \phi\sqrt{3}$ (Lemma~\ref{lem:phi-sqrt3}).
\item Enumerating all 1728 directed 3-edge paths $\vhost \to u_1 \to u_2
  \to u_3$ and computing the assembly~\eqref{eq:jk-assembly} with full
  9-shell vertex classification (the 4-shell coarsening $B_0..B_3$ that
  suffices at $k \leq 2$ does NOT suffice at $k=3$; cf. THEO-DSL-10 reasoning
  fragment 0618.md).
\item Decomposing the result in the orthonormal basis $\{\nrho, \nax,
  \nperpthree, \ndiff\}$ via direct dot products (the frame is exactly
  orthonormal in $\mathbb{Q}[\phi]^4$).
\item Extracting closed forms in $\sqrt{3}\cdot\mathbb{Q}[\phi]/3$ via mpmath/
  PSLQ at 80-digit precision in the basis $\{1, \phi, \sqrt{3}, \sqrt{3}\phi\}$.
\item Verifying the resulting closed-form expressions to machine-double
  precision in the verify script.
\item Cross-checking via the vertex-aligned $k=1, k=2, k=3$ targets
  ($3/\phi^2$, $-9/\phi^2$, $-126 + 81\phi$), the THEO-DSL-9 reproduction
  at $k=2$ face ($-7/\phi^2$, $2\phi - 6$), and the $V_4$-forced vanishing
  of the two non-invariant components (machine zero across all 30 faces).
\end{enumerate}
See \texttt{code/verify\_face\_alpha3\_closure.py} for the full verification
(all 30 PASS checks). The verbatim reasoning trail is captured in
\texttt{hardened\_theorems/reasoning/0620.md}.

\begin{antierasure}
The $\sqrt{3}$ factor in the THEO-DSL-11 closed forms is a substantive new
feature relative to the entire prior F.1 corpus, which has lived in
$\mathbb{Q}[\phi]$. Should a future multi-AI review surface either (i) a
counter-derivation in clean $\mathbb{Q}[\phi]$ (which would contradict
Theorem~\ref{thm:parity-structure} unless the $V_4$-invariant subspace
admits a basis other than $\{\nrho, \nax\}$ that hides the $\sqrt{3}$), or
(ii) a coefficient value different from those stated in
Theorem~\ref{thm:dsl11}, this artifact must be retained verbatim and a
separate revision/supersession artifact created, following the procedure
established at Patches 0615--0617 for the THEO-DSL-8 $\to$ THEO-DSL-9
correction. The author has not been able to construct a $\mathbb{Q}[\phi]$
basis for $\mathrm{span}\{\nrho, \nax\}$ that absorbs the $\sqrt{3}$
without introducing it elsewhere; the parity-dependent algebraic structure
of Theorem~\ref{thm:parity-structure} therefore appears to be a robust
feature of the face-aligned variant, but multi-AI verification is invited
on this point.

\paragraph{Why no $\mathbb{Q}[\phi]$-basis change can absorb the $\sqrt 3$ --- short proof sketch.}
The $\sqrt{3}$ in the closed forms is not a feature of the choice of frame
$\{\nrho, \nax\}$; it originates in the path-product weight
$\prod_{j=1}^{k}(\hat e_j \cdot \nperpF)$. Each factor
$(\hat e_j \cdot \nperpF)$ lies in $(1/\sqrt{3})\cdot\mathbb{Q}[\phi]$
(Corollary~\ref{cor:nFperp-sqrt3}), so the path product carries a
$\sqrt{3}^{-k}$ overall factor that is invariant under any change of
$\mathbb{Q}[\phi]$-basis of the ambient $\mathbb{Q}[\phi]^4$. If $B'$ is any
$\mathbb{Q}[\phi]$-basis of $\mathrm{span}\{\nrho, \nax\}$, expressing
$\jvec_k$ in $B'$ amounts to applying a $\mathbb{Q}[\phi]$-linear
transformation; this maps $\sqrt{3}^{-k}\cdot\mathbb{Q}[\phi]$ entries to
$\sqrt{3}^{-k}\cdot\mathbb{Q}[\phi]$ entries (since
$\mathbb{Q}[\phi]\cdot\sqrt{3}^{-k}\cdot\mathbb{Q}[\phi] =
\sqrt{3}^{-k}\cdot\mathbb{Q}[\phi]$ as $\mathbb{Q}[\phi]$-modules). At odd
$k$, $\sqrt{3}^{-k} \in \sqrt{3}^{-1}\cdot\mathbb{Q}[\phi]$, so the entries
remain in $\sqrt{3}^{-1}\cdot\mathbb{Q}[\phi] = (\sqrt 3 / 3)\cdot\mathbb{Q}[\phi]$,
i.e., outside $\mathbb{Q}[\phi]$. No $\mathbb{Q}[\phi]$-basis change can
remove the $\sqrt{3}$. This makes Theorem~\ref{thm:parity-structure}
basis-invariant: the parity-dependent algebraic-class membership is an
intrinsic feature of the face-aligned Reading C coefficients, not an
artifact of frame choice.
\end{antierasure}

% BEGIN GENERATED IDENTIFIER APPENDIX -- do not edit by hand
\clearpage
\section*{Appendix: Programme identifiers used in this paper}
\addcontentsline{toc}{section}{Appendix: Programme identifiers}

\noindent This programme tracks its results, assumptions and unresolved questions under short identifiers, so that a claim and its dependencies can be cited precisely. Prefixes denote the kind of item: \texttt{THEO-} a proved theorem, \texttt{FI-} a foundational input the derivation assumes, \texttt{PRED-} a registered prediction, \texttt{CONJ-} a conjecture, \texttt{OPEN-} an unresolved question, \texttt{PH-} a problem history. Those appearing in this paper are glossed below; no external document is needed to read them.

\begin{description}
  \item[\texttt{OPEN-FP-F1-1}] \hfill \\ Extend the F.1 substrate-locality umbrella theorem (Theorem 7.1) to \$\textbackslash\{\}mathcal\{O\}(\textbackslash\{\}delta\textasciicircum{}2)\$ by deriving the second-shell inner-product and edge-projection identities for the 600-cell, and determine whether the parallel-to-\$\textbackslash\{\}hat\{n\}\$ structure of \$\textbackslash\{\}vec\{J\}\_\{\textbackslash\{\}text\{DI-net\}\}(v\_\{\textbackslash\{\}text\{host\}\})\$ survives at second order. **Resolution**: both sub-questions closed — structural sub-question at publication-grade Layer 3 unconditional (\$\textbackslash\{\}vec\{j\}\_k(v\_\{\textbackslash\{\}text\{host\}\}) \textbackslash\{\}parallel \textbackslash\{\}hat\{n\}\$ at every \$k \textbackslash\{\}geq 1\$ via THEO-DSL-4) + coefficient sub-question at sketch-document Layer 3 under (H5) (\$\textbackslash\{\}alpha\_2 = -9/\textbackslash\{\}phi\textasciicircum{}2 = -9(2-\textbackslash\{\}phi) \textbackslash\{\}approx -3.438\$ via THEO-DSL-5 candidate; ratio \$\textbackslash\{\}alpha\_2/\textbackslash\{\}alpha\_1 = -3/2\$ exactly).
  \item[\texttt{OPEN-FP-F1-5}] \hfill \\ Extend or reformulate F.1 Theorems 5.1 (host-first-shell projection), 5.2 (first-shell perpendicularity), and 7.1 (substrate-locality umbrella) under edge-aligned Reading C (\$\textbackslash\{\}hat\{n\}\$ along a 600-cell short-edge direction; residual \$D\_5\$ symmetry at edge midpoint per Patch 0596 Remark 5.1 anti-erasure correction — see below) and face-aligned Reading C (\$\textbackslash\{\}hat\{n\}\$ perpendicular to a triangular face; residual \$C\_s\$ symmetry under vertex-host convention per Patch 0597 Remark 7.1 anti-erasure correction — see below).
  \item[\texttt{OPEN-SD-CHIR-PRIMITIVE}] \hfill \\ Derive the universe's primitive chirality bias from a single substrate-level mechanism (current leading candidate: primitive 4D direction \$\textbackslash\{\}hat\{n\}\$ aligned with a 600-cell host vertex per Reading C) that simultaneously accounts for the five observable manifestations across the CPP corpus: K3-doublet chirality (Capotauro \$\textbackslash\{\}Delta p\_\{LR\}\$, mass-mixing sector), W-bracelet V-A coupling (electroweak, massless helicity limit), electromagnetic handedness (E×B / Poynting / Maxwell right-hand rule and qDP/eDP polarization patterns), thermodynamic causal arrow (DI-bit propagation direction / PCD cycle ordering / Conscious Moment sequence), and cosmological vacuum asymmetry (matter-antimatter / baryogenesis / Capotauro nucleation-event sign-selection).
  \item[\texttt{THEO-DSL-1}] \hfill \\ Perturbation-Theory Propagation Rule + Shell-Locality at \$\textbackslash\{\}mathcal\{O\}(\textbackslash\{\}delta\textasciicircum{}1)\$ (**first F-line flagship theorem under THEO-DSL-N sub-prefix convention**; publication-grade Layer 3 **unconditional**; first of three F.1 v1.0 SHIPPED theorems closing OPEN-SD-CHIR-PRIMITIVE manifestation (iv) thermodynamic causal arrow at sketch-document Layer 3 via Theorem 7.1 substrate-locality umbrella)
  \item[\texttt{THEO-DSL-10}] \hfill \\ A Dynamical Substrate Law theorem confirmed by multi-model review.
  \item[\texttt{THEO-DSL-11}] \hfill \\ A Dynamical Substrate Law theorem with empirical validation carried out at fourth order.
  \item[\texttt{THEO-DSL-12}] \hfill \\ A Dynamical Substrate Law theorem in the registered series running from DSL-1 to DSL-12.
  \item[\texttt{THEO-DSL-2}] \hfill \\ First-Shell Geometric Identities at Vertex-Aligned Reading C in 600-Cell (**second F-line flagship theorem under THEO-DSL-N sub-prefix convention**; publication-grade Layer 3 **conditional on G1** first-shell inner-product primitive; combines two paper-body theorems of F.1 v1.0 SHIP into a single registry entry per F-line flagship convention)
  \item[\texttt{THEO-DSL-3}] \hfill \\ Substrate-Locality Umbrella Theorem (**third F-line flagship theorem under THEO-DSL-N sub-prefix convention**; **sketch-document Layer 3** — not independently hardened, assembled from THEO-DSL-1 + THEO-DSL-2 + icosahedral rank-1 sum identity; **closes OPEN-SD-CHIR-PRIMITIVE manifestation (iv) thermodynamic causal arrow** — fourth manifestation closure under the umbrella)
  \item[\texttt{THEO-DSL-4}] \hfill \\ \$\textbackslash\{\}mathcal\{O\}(\textbackslash\{\}delta\textasciicircum{}k)\$ Substrate-Current Parallel-to-\$\textbackslash\{\}hat\{n\}\$ Structural Theorem at All Orders \$k \textbackslash\{\}geq 1\$ (**fourth F-line flagship theorem under THEO-DSL-N sub-prefix convention**; publication-grade Layer 3 **unconditional** for the structural sub-question; **first Route C sequence 1 artifact** under the OPEN-FP-F1-1 trajectory; **closes the structural sub-question of OPEN-FP-F1-1** at all orders \$k \textbackslash\{\}geq 1\$ via symmetry-only argument; coefficient \$\textbackslash\{\}alpha\_k\$ at \$k \textbackslash\{\}geq 2\$ remains OPEN as target of Route C sequence 2)
  \item[\texttt{THEO-DSL-5}] \hfill \\ A Dynamical Substrate Law theorem giving a real coefficient in the substrate rate function at the host vertex.
  \item[\texttt{THEO-DSL-7}] \hfill \\ Umbrella registration covering first- and second-order edge-aligned coefficient content in the Dynamical Substrate Law.
  \item[\texttt{THEO-DSL-8}] \hfill \\ The cross-shell edge orbit closure under the reflection stabilizer of face-aligned Reading C.
  \item[\texttt{THEO-DSL-9}] \hfill \\ The candidate coefficient theorem for the face-aligned branch of the substrate law.
  \item[\texttt{THEO-DSL-N}] \hfill \\ (family) A proposed family registration for the shell-perpendicularity pattern that recurs at every perturbative order.
\end{description}
% END GENERATED IDENTIFIER APPENDIX

\end{document}
